\chapter{Future Directions}
\label{ch:future-directions}

\textit{This chapter outlines two tiers of future work: semester-length
projects (3--4 months) suitable for a graduate course or independent study,
and thesis-level research directions (6--12 months) that could form the
core of an MS thesis or PhD chapter.  Each direction includes a problem
statement, a suggested approach, and concrete deliverables or research
questions.}


% ═══════════════════════════════════════════════════════════════════════
\section{Semester Projects (3--4 months)}
\label{sec:future-semester}

These projects are scoped to produce a concrete artifact and a short
evaluation within one academic semester.


% ─────────────────────────────────────────────────────────────────────
\subsection{Real Conversation Evaluation}
\label{sec:future-real-eval}

\noindent\textbf{Problem.}
All current evaluation data is synthetic.  Metrics may not generalize to
real developer--AI conversations (see
Section~\ref{sec:lim-synthetic}).

\medskip\noindent\textbf{Approach.}

\begin{enumerate}
    \item Identify open-source projects that publish AI coding assistant
          logs (e.g.\ projects using Claude Code with public session logs,
          or GitHub Copilot Chat histories shared in issues).
    \item Collect 50--100 real conversations and anonymize them: replace
          developer names, project names, and any PII.
    \item Annotate entity mentions with ground-truth canonical names
          (two independent annotators, inter-annotator agreement measured
          with Cohen's $\kappa$).
    \item Run \sysname{}'s pipeline on the real data and compare B-cubed
          metrics against the synthetic results.
\end{enumerate}

\medskip\noindent\textbf{Deliverables.}

\begin{itemize}
    \item An anonymized real-conversation evaluation corpus.
    \item B-cubed precision/recall/F1 on real data vs.\ synthetic data.
    \item Analysis of failure modes unique to real conversations.
\end{itemize}


% ─────────────────────────────────────────────────────────────────────
\subsection{Stronger Baselines}
\label{sec:future-baselines}

\noindent\textbf{Problem.}
The current evaluation lacks strong baselines (see
Section~\ref{sec:lim-weak-baselines}).

\medskip\noindent\textbf{Approach.}

\begin{enumerate}
    \item Implement a \textbf{BERT entity matching} baseline: fine-tune a
          pre-trained BERT model on the training split to predict whether
          two entity mentions refer to the same canonical entity.
    \item Implement a \textbf{TF-IDF cosine similarity} baseline using
          character n-grams (see Exercise~\ref{sec:ex-new-baseline}).
    \item Optionally implement or adapt an existing entity linking system
          (e.g.\ BLINK, EntQA) to the technology-entity domain.
    \item Evaluate all systems on the same held-out test set with B-cubed
          metrics and McNemar's test for statistical significance.
\end{enumerate}

\medskip\noindent\textbf{Deliverables.}

\begin{itemize}
    \item At least two baseline implementations with reproducible
          evaluation scripts.
    \item A comparison table showing precision, recall, F1, and
          statistical significance for each system.
    \item Analysis of which mention types (abbreviations, misspellings,
          etc.) each system handles well or poorly.
\end{itemize}


% ─────────────────────────────────────────────────────────────────────
\subsection{User Study}
\label{sec:future-user-study}

\noindent\textbf{Problem.}
The practical value of a decision knowledge graph is unvalidated (see
Section~\ref{sec:lim-no-user-study}).

\medskip\noindent\textbf{Approach.}

\begin{enumerate}
    \item Design a controlled experiment with two groups of developers
          (minimum 10 per group).
    \item Both groups receive the same set of 5 decision-making tasks
          (e.g.\ ``Choose a caching strategy for this microservice given
          the existing architecture'').
    \item The treatment group has access to \sysname{} with a
          pre-populated knowledge graph.  The control group has access to
          the same information in a flat document (e.g.\ a README).
    \item Measure: decision consistency (do team members make the same
          choice?), time to decision, and subjective usefulness (Likert
          scale survey).
\end{enumerate}

\medskip\noindent\textbf{Deliverables.}

\begin{itemize}
    \item IRB-approved study protocol.
    \item Quantitative comparison of decision quality and time between
          groups.
    \item Qualitative analysis of how developers use the knowledge graph
          (or do not use it).
\end{itemize}


% ─────────────────────────────────────────────────────────────────────
\subsection{Dormant Stage Analysis}
\label{sec:future-dormant}

\noindent\textbf{Problem.}
Two pipeline stages (alias search, embedding similarity) rarely trigger
on current data (see Section~\ref{sec:lim-dormant-stages}).

\medskip\noindent\textbf{Approach.}

\begin{enumerate}
    \item Create an adversarial test set specifically designed to bypass
          the first four stages: mentions with creative abbreviations,
          domain-specific jargon, and ambiguous references.
    \item Measure how often stages 4 and 6 activate on this harder data.
    \item If they activate and resolve correctly, their value is
          demonstrated.
    \item If they do not, propose a simplified pipeline that removes them,
          and compare accuracy with and without.
\end{enumerate}

\medskip\noindent\textbf{Deliverables.}

\begin{itemize}
    \item An adversarial evaluation dataset (at least 100 hard mentions).
    \item Per-stage activation rates on easy vs.\ hard data.
    \item A recommendation: keep, modify, or remove the dormant stages.
\end{itemize}


% ─────────────────────────────────────────────────────────────────────
\subsection{Confidence Calibration}
\label{sec:future-calibration}

\noindent\textbf{Problem.}
Entity resolution confidence scores (avg 0.945) may not be well
calibrated---a score of 0.95 may not mean the prediction is correct 95\%
of the time.

\medskip\noindent\textbf{Approach.}

\begin{enumerate}
    \item Compute ECE on the current evaluation data using
          \codefile{evaluation/compute\_calibration.py}.
    \item If ECE is high, train a post-hoc calibration layer (e.g.\
          Platt scaling or isotonic regression) on the training split.
    \item Evaluate calibrated confidence scores on the test split.
    \item Visualize reliability diagrams (expected accuracy vs.\ predicted
          confidence per bin).
\end{enumerate}

\medskip\noindent\textbf{Deliverables.}

\begin{itemize}
    \item ECE before and after calibration.
    \item A calibration model that can be integrated into the pipeline.
    \item Reliability diagrams showing calibration improvement.
\end{itemize}


% ═══════════════════════════════════════════════════════════════════════
\section{Thesis-Level Research (6--12 months)}
\label{sec:future-thesis}

These directions are open-ended research problems that require sustained
investigation and could produce publishable results.


% ─────────────────────────────────────────────────────────────────────
\subsection{Multi-Agent Knowledge Sharing via MCP}
\label{sec:future-multi-agent}

\noindent\textbf{Research question.}
Can multiple AI coding agents share a decision knowledge graph through the
Model Context Protocol (MCP) to improve collective decision quality?

\medskip\noindent\textbf{Hypothesis.}
When multiple agents working on the same codebase share decisions through
a common knowledge graph, they will make more consistent decisions and
avoid contradictions compared to agents working in isolation.

\medskip\noindent\textbf{Methodology.}

\begin{enumerate}
    \item Deploy \sysname{} as an MCP server accessible to multiple AI
          agent instances.
    \item Design a multi-agent scenario: 3--5 agents working on different
          modules of the same project, each making technology decisions.
    \item Measure decision consistency (do agents choose compatible
          technologies?) with and without the shared knowledge graph.
    \item Analyze the knowledge graph for contradiction patterns
          (\ic{CONTRADICTS} edges) and resolution strategies.
\end{enumerate}


% ─────────────────────────────────────────────────────────────────────
\subsection{Cross-Project Knowledge Transfer}
\label{sec:future-cross-project}

\noindent\textbf{Research question.}
Can entity alignment across project knowledge graphs enable useful
knowledge transfer?

\medskip\noindent\textbf{Hypothesis.}
Decisions about a technology in one project (e.g.\ ``we moved from REST
to GraphQL because of N+1 query problems'') are relevant to similar
decisions in other projects using the same technology stack.

\medskip\noindent\textbf{Methodology.}

\begin{enumerate}
    \item Populate \sysname{} with data from at least 5 different projects
          that share some technology overlap.
    \item Implement entity alignment: match entities with the same
          canonical name across projects.
    \item Build a cross-project query interface (e.g.\ ``What has our
          organization decided about PostgreSQL?'').
    \item Evaluate whether cross-project context improves GraphRAG answer
          quality compared to single-project context.
\end{enumerate}


% ─────────────────────────────────────────────────────────────────────
\subsection{Real-Time MCP Capture}
\label{sec:future-realtime}

\noindent\textbf{Research question.}
Is real-time decision capture during AI-assisted development more
effective than post-hoc extraction from conversation logs?

\medskip\noindent\textbf{Hypothesis.}
Real-time capture preserves contextual signals (developer intent, code
state at decision time) that are lost in post-hoc analysis, leading to
higher-quality decision records.

\medskip\noindent\textbf{Methodology.}

\begin{enumerate}
    \item Implement a real-time capture mode in the MCP server that
          intercepts decisions as they happen during coding sessions.
    \item Compare decision quality (completeness of trigger, context,
          options, rationale) between real-time and post-hoc extraction.
    \item Measure latency impact: does real-time capture add perceptible
          delay to the coding workflow?
    \item Survey developers on the intrusiveness of real-time capture.
\end{enumerate}


% ─────────────────────────────────────────────────────────────────────
\subsection{Privacy-Preserving Federated Graphs}
\label{sec:future-federated}

\noindent\textbf{Research question.}
Can organizations share decision knowledge without exposing proprietary
implementation details?

\medskip\noindent\textbf{Hypothesis.}
A federated learning approach---where each organization trains a local
entity resolution model and shares only model updates (not raw data)---can
produce a globally improved resolution model without compromising privacy.

\medskip\noindent\textbf{Methodology.}

\begin{enumerate}
    \item Partition the evaluation data into 3--5 ``organizations'' with
          non-overlapping domain distributions.
    \item Implement federated averaging for the embedding-based resolution
          stage: each organization trains locally, then shares embedding
          model weight updates.
    \item Compare resolution quality of the federated model against
          locally-trained models and a centralized model (upper bound).
    \item Analyze privacy guarantees: what information can be inferred
          from shared model updates?
\end{enumerate}


% ─────────────────────────────────────────────────────────────────────
\subsection{Temporal Graph Reasoning}
\label{sec:future-temporal}

\noindent\textbf{Research question.}
Can temporal reasoning over the decision graph answer questions about
decision evolution, such as ``When did we stop using technology X?'' or
``How has our architecture philosophy changed over the past year?''

\medskip\noindent\textbf{Hypothesis.}
Adding temporal annotations (creation time, supersession time) to graph
edges enables a class of useful queries that static graphs cannot support.

\medskip\noindent\textbf{Methodology.}

\begin{enumerate}
    \item Extend the graph schema with temporal properties on all edges
          (valid-from, valid-to, creation timestamp).
    \item Implement temporal Cypher query patterns: point-in-time
          snapshots, interval queries, and trend analysis.
    \item Define a benchmark of 20 temporal questions and evaluate whether
          \sysname{} can answer them correctly.
    \item Compare against a baseline that treats the graph as static (no
          temporal reasoning).
\end{enumerate}


% ─────────────────────────────────────────────────────────────────────
\subsection{Scale and Performance}
\label{sec:future-scale}

\noindent\textbf{Research question.}
How does the system perform at 10$\times$, 100$\times$, and 1{,}000$\times$
the current scale (10K+ entities, 100K+ relationships)?

\medskip\noindent\textbf{Hypothesis.}
The current linear-scan entity resolution and exact nearest-neighbor
embedding search will become bottlenecks at scale, but can be replaced
with approximate algorithms without significant accuracy loss.

\medskip\noindent\textbf{Methodology.}

\begin{enumerate}
    \item Generate synthetic data at 10$\times$ (2,000 conversations),
          100$\times$ (20,000), and 1{,}000$\times$ (200,000) the current
          dataset size.
    \item Profile the full pipeline at each scale: entity resolution
          latency, graph query time, GraphRAG retrieval time, and memory
          usage.
    \item Replace exact nearest-neighbor search with ANN indexes
          (FAISS HNSW or pgvector IVFFlat) and measure the
          accuracy--latency trade-off.
    \item Implement graph sampling for large-neighborhood traversals in
          the GraphRAG pipeline.
    \item Report scaling curves: latency and accuracy as functions of
          dataset size.
\end{enumerate}
